\chapter{When: Retrieval-Efficient Acquisition Selection}
\label{ch:when}

\basedon{\fullcite{ulker2027acquisitions}}

\section{Introduction}

\section{Problem formulation}

\section{Method}
\subsection{Any-subset classifier}
\subsection{Per-region acquisition calendar (greedy search)}
\subsection{Label-free learned stopping head}

\section{Data and experimental protocol}

\section{Results}
\subsection{Searched calendars versus equal spacing}
\subsection{Learned stopping}
\subsection{Comparison and combination with ELECTS}
\subsection{Transfer across years}

\section{Why oracle-guided acquisition policies fail}
\label{sec:when:oracle}

\section{Per-field adaptive date choice}

\section{Discussion and limitations}

\section{Conclusion}

\begin{bridge}
Chapters~\ref{ch:which} and~\ref{ch:when} reduce the input along two
independent axes. If both reductions hold at once, the data volume per field
shrinks by their product. The next chapter tests whether they compose, and
how to choose bands and dates together under one data-volume budget.
\end{bridge}
